% PINN revision text. Figure numbering is resolved through \label/\ref.

\subsubsection{Naturalistic replay and recording-disjoint construct validity}
The field-query interface was examined on recorded traffic before its use in
closed-loop decision making. At each sensing instant, the measured vehicle
states are transformed to an ego-local frame and encoded as source, transport,
diffusion, occupancy, and road-support channels. The resulting local field is
then mapped back to the recorded world pose for visualization. In
Fig.~\ref{fig:pinn-naturalistic-overlay}, this mapping follows the official
dataset calibration: metric world coordinates are converted to orthophoto
pixels using the recording-specific scale, and all vehicles are drawn from
their recorded bounding boxes at the same frame. Thus, the heatmap, roadway,
and traffic participants share one coordinate system rather than being
juxtaposed as independent local-grid images.

\begin{figure*}[t]
  \centering
  \includegraphics[width=\textwidth]{pinn_naturalistic_orthophoto_overlay.pdf}
  \caption{Numerical and PINN risk-field queries replayed on naturalistic
  traffic recordings. Each row shows the same recorded frame and vehicle
  configuration under the prospective numerical field (left) and the
  context-conditioned PINN query (right). The displayed frame is selected by
  maximum integrated numerical risk among frames with a passenger or heavy
  road-vehicle ego, and both columns use a common scene-wise normalization.
  The inD example is recording-disjoint; rounD and exiD are external
  qualitative replays and are not pooled into the calibrated validation
  statistics. Red, orange, and light-blue boxes identify the ego, heavy, and
  surrounding vehicles, respectively.}
  \label{fig:pinn-naturalistic-overlay}
\end{figure*}

\paragraph{Role of the context-conditioned PINN.}
Let $\mathcal{S}_{H}$ denote the numerical transport--diffusion operator over
the prospective horizon $H$ and let $\mathbf{q}_{t}$ collect only information
available at sensing time $t$, including the instantaneous source,
longitudinal and lateral transport, diffusion, occupancy, and road-support
channels. The numerical teacher evaluates
$R^{\mathrm{num}}_{t}=\mathcal{S}_{H}(\mathbf{q}_{t})$ by repeated integration,
whereas the frozen network evaluates
$\widehat R_{t}=f_{\theta}(\mathbf{q}_{t})$ in one differentiable computation.
The purpose of $f_{\theta}$ is therefore to amortize a nonlocal finite-horizon
operator across scenes, not to relabel the teacher as empirical ground truth.
Its fixed computation graph supports batched GPU queries and direct spatial
derivatives, while the physics residual and gradient-matching losses restrict
solutions that fit snapshot values but violate the teacher dynamics. The
network was calibrated on inD recordings 01--05 and Highway-env seeds 0--2
from four topologies; inD recordings 06--11 and Highway-env seeds 100--102
were excluded from parameter fitting and normalization. Future trajectories,
collision outcomes, RL returns, and social scores were not used as PINN
training targets.

This context-conditioned formulation replaces the earlier recording-specific
surrogate with one operator that is queried from the current scene tensor. The
recordings supply heterogeneous traffic states and road geometries at which
the operator is learned; they do not supply a hindsight risk label to the
training loss. Consequently, naturalistic behavior enters through the measured
state and transport context, while the transport--diffusion equation remains
the mechanistic prior. The numerical solver is retained as the auditable
reference, whereas the PINN provides a fixed-depth, differentiable, and
batchable implementation for repeated online queries.

The overlay in Fig.~\ref{fig:pinn-naturalistic-overlay} is an interpretability
and coordinate-registration check, not an accuracy test. In particular, a
smoother or smaller heatmap does not establish a safer or more naturalistic
planner. We consequently test the prospective spatial construct against the
observed three-second future occupancy envelope on recording-disjoint inD
data. This label is formed only from subsequently observed vehicle trajectories
and is independent of both the PDE source and the RL reward. The instantaneous
source obtains AUROC/AUPRC values of $0.8249/0.4167$, propagation by the numerical
teacher raises them to $0.9096/0.5836$, and the frozen PINN obtains
$0.8657/0.5143$. Relative to the source-only control, the PINN gains are
$0.0408$ (recording-bootstrap $95\%$ CI $[0.0190,0.0576]$) and $0.0976$
($[0.0720,0.1207]$), respectively. Thus, the learned query retains predictive
structure associated with the nonlocal propagation rather than merely
reproducing the instantaneous source amplitude. The numerical teacher remains
the higher-fidelity reference, and neither field is interpreted as a
calibrated collision probability.

\subsubsection{Decision-level numerical-to-PINN substitution}
The influence of approximation error is evaluated at three successive
interfaces. First, held-out fields are compared only where the numerical
gradient is active, because an angle at a near-zero gradient is undefined and
cannot direct a maneuver. Second, synchronized simulator states are passed
through both backends. The resulting eight-dimensional descriptor
$\mathbf{o}^{R}_{t}=[R_{0},R_{5},R_{10},R_{20},\partial_xR,
\partial_yR,R_{L},R_{R}]$ is normalized by fixed calibration scales and clipped
to $[-3,3]$ before it enters the frozen policy. This test measures
$\|\Delta\mathbf{o}^{R}_{t}\|$ and
$\|\pi_{\phi}(\mathbf{o}_{t},\mathbf{o}^{R,\mathrm{PINN}}_{t})-
\pi_{\phi}(\mathbf{o}_{t},\mathbf{o}^{R,\mathrm{num}}_{t})\|_{2}$ directly;
it does not assume that a percentage field-gradient error produces the same
percentage action error. Third, for each traffic seed, the same policy is
executed in paired closed-loop episodes while replacing only the query
backend. Collision, minimum time to collision, corridor-risk exposure,
progress, speed, jerk, imposed follower deceleration, backward risk flux, and
latency are then compared. Confidence intervals are obtained by paired
bootstrap over traffic seeds, and multiplicity is controlled across the
predeclared scenario--metric comparisons.

Equivalently, the analysis follows the composed map
$R\xrightarrow{\mathcal D}\mathbf{o}^{R}
\xrightarrow{\pi_{\phi}}\mathbf{u}\xrightarrow{\mathcal E}\mathbf{Y}$,
where $\mathcal D$ is the descriptor extractor, $\pi_{\phi}$ is the frozen
policy, $\mathcal E$ denotes closed-loop environment evolution, and
$\mathbf{Y}$ collects safety, progress, comfort, and externality outcomes.
For an individual synchronized state, a local relation
$\|\Delta\mathbf{u}_{t}\|_{2}\leq
L_{\pi,t}\|\Delta\mathbf{o}^{R}_{t}\|_{2}$ may be written, but no global
Lipschitz constant is assumed for the neural policy. Instead, both terms are
measured on the paired state sequence and the residual effect is tested by the
backend intervention at the level of $\mathbf{Y}$.

The completed merge swap is a non-inferiority-style substitution test, not
evidence that the PINN is intrinsically less conservative. Across ten paired
seeds, the descriptor RMSE is $4.72\times10^{-3}$, the mean action-vector
difference is $6.68\times10^{-5}$, and no synchronized query exceeds the
predeclared action tolerance $\|\Delta\mathbf{u}\|_{2}=0.05$. Both backends
produce zero collisions and equal episode lengths. Mean progress changes from
$585.533$ to $585.473$~m ($\Delta=-0.061$~m; $95\%$ paired-bootstrap CI
$[-0.076,-0.046]$~m), mean speed changes by $-0.0074$~m/s, and mean absolute
jerk changes by $-2.52\times10^{-4}$~m/s$^{3}$. These magnitudes support
behavior-preserving substitution for this merge policy, but not superior
planning. The isolated CUDA field kernel is faster than the CPU numerical
core ($2.04$ versus $2.64$~ms), whereas the complete online query remains
slightly slower after context construction and host--device transfer
($5.80$ versus $5.21$~ms). The deployment advantage is therefore a
differentiable, batchable, fixed-cost surrogate whose present end-to-end
latency still requires on-device context construction; the numerical solver
remains available as the selectable validation reference and fallback.

These experiments separate three claims that were conflated in the original
manuscript. First, propagation is evaluated against an instantaneous-source
control using future occupancy that is unavailable to either field model;
this tests whether the nonlocal construction carries prospective interaction
information. Second, the numerical-to-PINN swap tests whether amortization
preserves the descriptors and decisions of the physics-based reference. Third,
the downstream observation--reward factorial evaluates whether supplying a
field descriptor, changing the objective, or their interaction alters
closed-loop driving. Accordingly, improvements over native RL are attributed
to the field representation and the social objective only when supported by
that factorial comparison. The demonstrated advantage of the present PINN over
online numerical integration is its differentiable, batchable query interface
with behaviorally equivalent merge execution; the evidence does not support a
claim that smoothing the numerical field by itself improves planning. This
separation retains the mechanistic prior supplied by the PDE while preventing
surrogate approximation, risk validity, and policy benefit from being inferred
from one reconstruction statistic. Thus, relative to the instantaneous source,
the propagated representation carries additional recording-disjoint
information about future occupancy; relative to online numerical integration,
the PINN retains policy-relevant behavior through an amortized query; and any
driving improvement over native RL is attributed only through the separate
observation--reward factorial.

% ---------------------------------------------------------------------------
% Response letter text
% ---------------------------------------------------------------------------

\noindent\textbf{Reviewer comment:}
\emph{How can the authors demonstrate that the learned field captures genuine
interaction risk rather than reproducing manually injected surrogate
variables?}

\noindent\textbf{Response:}
We agree with the distinction underlying this comment. Reconstructing a
hand-designed teacher establishes surrogate fidelity, but it cannot establish
that the teacher variable corresponds to interaction risk. We therefore
separated the revised validation into two hypotheses and ensured that the
second hypothesis uses evidence unavailable during PINN training.

For surrogate fidelity, the network is trained only to approximate the causal
three-second transport--diffusion operator and its spatial gradient, subject
to a physics residual. The inputs contain the scene state available at the
current sensing instant; neither future trajectories nor collision, reward,
or social-score labels enter the teacher cache or PINN loss. Calibration uses
inD recordings 01--05 and Highway-env training seeds 0--2. Domain-specific
normalization is fitted on that calibration set only and then frozen.

For construct validity, we evaluate on the disjoint inD recordings 06--11 and
form an independent target from the vehicle occupancy observed during the
subsequent three seconds. This target is not a PDE source component and is not
used by the RL objective. We compare three progressively richer quantities:
the instantaneous hand-designed source, the propagated numerical field, and
the frozen PINN query. The source-only control obtains AUROC/AUPRC values of
$0.8249/0.4167$; propagation raises these values to $0.9096/0.5836$; and the
PINN retains $0.8657/0.5143$. Relative to source only, the PINN gains are
$0.0408$ (recording-bootstrap 95\% CI $[0.0190,0.0576]$) and $0.0976$
($[0.0720,0.1207]$). Because the resampling unit is the recording rather than
the frame, the intervals do not treat temporally correlated snapshots as
independent samples. The result shows that the propagated representation
contains recording-disjoint information about where traffic will occupy the
scene beyond the instantaneous injected source. It does not prove causality,
nor do we describe either output as a calibrated probability of collision.

This comparison is a direct falsification test of the ``source-copying''
explanation. If the learned field contained only a monotone reproduction of
the current source amplitude, it could not improve the ranking of a target
constructed exclusively from later recorded occupancy relative to that source
under a recording-disjoint split. The positive paired gains therefore reject
that restricted explanation for the tested construct, while the remaining gap
to the numerical teacher quantifies information lost by amortization. We have
made this evidential boundary explicit: the experiment supports prospective
interaction relevance, not a causal or probabilistic definition of risk and
not human-preference learning.

Finally, Fig.~\ref{fig:pinn-naturalistic-overlay} now uses the official
world-to-orthophoto calibration and recorded bounding boxes at a common frame.
The inD row is part of the recording-disjoint test; rounD and exiD are clearly
labeled zero-shot qualitative replays and are not pooled into the numerical
claim. This visual check verifies registration and qualitative topology, while
the future-occupancy test supplies the independent quantitative evidence.

\noindent\textbf{Changes in the manuscript:}
We added the subsection ``Naturalistic replay and recording-disjoint
construct validity,'' specified the calibration/evaluation recordings and
label construction, reported source-only, numerical, and PINN results with
recording-bootstrap intervals, and narrowed the claim from ``learned collision
risk'' to ``predictive relevance for near-future interaction occupancy.'' We
also replaced the local-grid montage with orthophoto-calibrated replays and
removed error statistics from the qualitative figure.

\medskip
\noindent\textbf{Reviewer comment:}
\emph{The reported gradient error may systematically mislead acceleration,
deceleration, or lane-change decisions, but the paper does not analyze how PINN
approximation error propagates into policy performance.}

\noindent\textbf{Response:}
We agree that field-level gradient error alone is insufficient. We also note
that an angular error is not a direct percentage error in steering or
acceleration: the angle is unstable where the reference gradient is nearly
zero, and the policy receives eight normalized descriptors together with the
native traffic observation rather than applying the raw gradient as a control
command. We therefore restrict the angular diagnostic to active-gradient
cells and test error propagation at the actual policy interface.

The revised analysis has three levels. First, it reports topology-specific
field correlation, active-gradient angle, and equal-mass hotspot overlap on
held-out Highway-env frames. The aggregate active-gradient error is
$13.8^{\circ}$, with $7.4^{\circ}$ in merge; roundabout is explicitly reported
as the weakest topology rather than hidden in an average. Second, for
synchronized states, both backends produce the eight-dimensional vector
$[R_{0},R_{5},R_{10},R_{20},\partial_xR,\partial_yR,R_L,R_R]$. The risk and
gradient channels are normalized using fixed calibration scales and clipped
to $[-3,3]$ before the same frozen policy is queried. On ten merge seeds, the
descriptor RMSE is $4.72\times10^{-3}$, the mean action-vector difference is
$6.68\times10^{-5}$, and none of the paired actions exceeds the predeclared
tolerance $\lVert\Delta\mathbf{u}\rVert_2=0.05$. This directly measures the
policy's local sensitivity and avoids assuming an unavailable global
Lipschitz constant.

Formally, we audit the entire composition
$R\xrightarrow{\mathcal D}\mathbf{o}^{R}
\xrightarrow{\pi_{\phi}}\mathbf{u}\xrightarrow{\mathcal E}\mathbf{Y}$ rather
than treating an angular field statistic as a control error. Although a local
bound $\|\Delta\mathbf{u}_{t}\|_{2}\leq
L_{\pi,t}\|\Delta\mathbf{o}^{R}_{t}\|_{2}$ exists wherever the policy is
locally Lipschitz, we do not estimate or assume a global constant. Paired
queries at identical states directly measure descriptor and action changes,
and the backend swap then measures whether their accumulation alters safety,
progress, comfort, or externality outcomes.

Third, we perform an actual closed-loop intervention: the simulator, policy,
traffic seed, and reward are fixed, and only the numerical/PINN query backend
is exchanged. Both merge backends complete all ten episodes without collision
and with identical episode lengths. PINN-minus-numerical progress is
$-0.061$~m (95\% paired-bootstrap CI $[-0.076,-0.046]$~m) over approximately
$586$~m, while mean speed and mean absolute jerk change by only
$-0.0074$~m/s and $-2.52\times10^{-4}$~m/s$^3$, respectively. We interpret
this as behavior-preserving substitution for the tested merge policy, not as
evidence of universal robustness or superior planning. The claim for any
additional topology is gated by the same paired descriptor, action, safety,
progress, comfort, and externality checks.

We also revised the computational claim. A paired batch-one benchmark uses 20
held-out merge frames, ten warm-up calls per backend, and excludes model
loading and rendering. The isolated CUDA network kernel is faster than the CPU
numerical field core ($2.04$ versus $2.64$~ms), but the complete CUDA query is
currently $5.80$~ms versus $5.21$~ms because scene encoding and host--device
transfer remain off-device. Policy inference, field query, environment step,
and rendering are reported as separate quantities. Thus, the present benefit
is a differentiable and batchable amortized operator with verified behavioral
interchangeability in merge; end-to-end acceleration is not claimed until
context construction is moved on-device.

\noindent\textbf{Changes in the manuscript:}
We added the subsection ``Decision-level numerical-to-PINN substitution,''
defined the active-gradient mask and descriptor preprocessing, and reported
the synchronized descriptor/action test together with the paired closed-loop
merge result. The runtime table now separates field-core latency,
context/transfer overhead, and complete query latency. The text presents the
merge result as a bounded non-inferiority result and explicitly reports the
weak roundabout fidelity instead of extrapolating from the aggregate error.
